%------------------------------------------------------------------------------
% \Introduction
%------------------------------------------------------------------------------
\chapter*{\vspace*{-\baselineskip}\begin{center}ВВЕДЕНИЕ\end{center}\vspace*{-\baselineskip*2}}
\addcontentsline{toc}{chapter}{\hspace*{-1em}Введение}\par

%------------------------------------------------------------------------------
% Text
%------------------------------------------------------------------------------
\hspace*{12.5 mm}В последние годы значительный интерес вызывает аппаратная реализация нейронных 
сетей, поскольку программные решения, работающие на центральных процессорах 
(CPU) и графических процессорах (GPU), не всегда обеспечивают требуемую 
скорость обработки и энергоэффективность. 

В связи с этим использование 
специализированных аппаратных ускорителей на базе FPGA (Field-Programmable 
Gate Array) или ASIC (Application-Specific Integrated Circuit) становится 
актуальной задачей. Аппаратная реализация нейронной сети в виде IP-ядра 
позволяет значительно повысить производительность и снизить задержки обработки 
данных, что особенно важно для встроенных систем.

При реализации нейронных сетей на базе CPU и GPU как правило используются 
стандартизированные типы данных. При реализации на базе FPGA появляется 
возможность использовать для представления параметров нейронных сетей типов 
данных, обеспечивающих различную точность. Причем выбор точности представления 
напрямую будет влиять на аппаратные затраты\cite{ITS2024}. 

Целью данного проекта является разработка IP-ядра нейронной сети 
прямого распространения для распознавания рукописных цифр. 
Необходимо не только создать высокоэффективное аппаратное решение, но и в 
продемонстрировать его практической применимости в условиях ограниченных 
вычислительных ресурсов. Разработанное IP-ядро может быть использовано в 
системах реального времени, таких как интеллектуальные камеры, портативные 
устройства и робототехнические комплексы, где требуется автономная и быстрая 
обработка визуальной информации без привлечения внешней вычислительной 
инфраструктуры.

В данном проекте 
используется обученное двумерное разделяемое преобразование (LST2D), которое 
рассматривается как новый тип слоя нейронной сети. Он следует концепции 
распределения веса. 

